\section{Técnicas e instrumentos, validación y confiabilidad}

Las técnicas principales serán el análisis documental de fuentes abiertas, el
procesamiento de series temporales, el entrenamiento comparativo de modelos y la
evaluación funcional de los módulos de apoyo a la decisión.

\begin{table}[H]
\centering
\caption{Técnicas e instrumentos de recolección y análisis}
\begin{tabular}{p{4cm}p{5cm}p{5cm}}
\toprule
\textbf{Técnica} & \textbf{Instrumento} & \textbf{Fuente / aplicación} \\
\midrule
Análisis documental & Ficha de extracción de datos & UNData, Kaggle y fuentes nacionales de exportación. \\
Procesamiento de datos & Scripts en Python & Limpieza, normalización, rezagos y variables temporales. \\
Modelado predictivo & SARIMA, LSTM y SVR & Entrenamiento y comparación de modelos. \\
Evaluación de modelos & MAE, RMSE, MAPE & Medición del error de pronóstico. \\
Validación funcional & Cuestionario de utilidad percibida & Productores de palta de La Libertad. \\
\bottomrule
\end{tabular}
\end{table}

\textbf{Validación y confiabilidad:} la calidad de datos será verificada mediante
revisión de valores faltantes, duplicados, rangos atípicos y consistencia
temporal. La confiabilidad del modelo se evaluará con partición temporal,
comparación de métricas y, de ser posible, validación cruzada para series de
tiempo. La utilidad del sistema se validará con productores mediante criterios
de claridad, pertinencia y facilidad de interpretación.
